\chapter{Improving the slope to $C^2$}
\label{chap:smooth-grains}

The Lipschitz slope supplied by the grain proposition has much more
structure than an arbitrary Lipschitz function. The local grains and
the global grains constrain the same shaded set. A projection theorem
turns these simultaneous constraints into affine approximation on
many heights. Repeating the argument at a finite sequence of scales
then permits smoothing without losing the global grain bound.

\begin{proposition}[$C^2$ grains; WZ2, Proposition 6.4]
\label{prop:c2-grains}
\lean{Kakeya.Assouad.PureWZ2C2GrainsFromCriticalStatement,Kakeya.Assouad.pure_wz2_node05_c2_grains}
\leanok
\uses{def:extremal,def:plane-map,def:ad-set}
For every $\eps,\delta_0>0$, there is an $\eps$-extremal
configuration at some $0<\delta\leq\delta_0$ whose shaded union $E$
has the following properties. There is $f\in C^2([-1,1])$ with
$\|f\|_{C^2}\leq1$ such that
\[
              (E\cap\{z=z_0\})\cdot(1,f(z_0),0)
\]
is a $(\delta,1-\sigma,\delta^{-\eps})_1$-ADset for every $z_0$.
The same configuration has a $1$-Lipschitz plane map $V:E\to S^2$
with the local grain property of Proposition~\ref{prop:lipschitz-grains}.
\end{proposition}
\begin{proof}
\uses{prop:lipschitz-grains,lem:locally-affine,lem:affine-hierarchy,
  cor:trapezoid-hierarchy,prop:trapezoid-smoothing,lem:mild-rescaling}
\leanok
Apply the affine approximation argument successively to the global
slope. Retain one common refinement of the shading, so that the local
grain estimates remain available at every step. The approximating
graphs lie in nested, horizontally separated trapezoids. Smooth their
successive affine corrections using disjointly supported cutoffs.
The derivative estimates follow by summing the correction bounds over
the finite hierarchy. The smoothed function agrees with the original
slope to within the fine scale on every retained height. This changes
each global projection by at most a fixed multiple of $\delta$.
Finally, a mild rescaling normalizes the $C^2$ and Lipschitz bounds.
\end{proof}

\begin{proposition}[The WZ1 smooth grain statement; WZ1, Proposition 5.1]
\label{prop:wz1-c2-grains}
\lean{Kakeya.Assouad.WZ1C2GrainPackageStatement}
\uses{def:wz1-extremal,def:plane-map,def:ad-set}
With WZ1 extremality in place of WZ2 extremality, the conclusion of
Proposition~\ref{prop:c2-grains} holds at arbitrarily small scales:
the global grain slope has $C^2$ norm at most $1$, and the same
shaded union retains its local grains and $1$-Lipschitz plane map.
\end{proposition}
\begin{proof}
\uses{prop:wz1-grains,cor:trapezoid-hierarchy,prop:trapezoid-smoothing}
This is the conclusion of WZ1 Section 5. WZ2 uses the same smoothing
argument after replacing the source of the Lipschitz grains.
\end{proof}

\section{The projection alternative}

\begin{theorem}[A projection dichotomy; WZ1, Theorem 5.2, corrected orientation]
\label{thm:wz1-projection}
\lean{Kakeya.Assouad.WZ1ProjectionDichotomyConclusion}
\uses{def:ad-set}
For every $\eps>0$, there are $\eta,\delta_0>0$ such that the
following holds at $0<\delta\leq\delta_0$. Let
$F,G\subset[0,1]^2$ be $(\delta,1,\delta^{-\eta})_2$-ADsets.
At least one of the following alternatives holds:
\begin{enumerate}
\item There is an affine line $\ell$ and its perpendicular line
$\ell^\perp$ through the origin such that
\[
 \cE_\delta(F\cap N_\delta(\ell^\perp))\geq\delta^{\eps-1},
 \qquad
 \cE_\delta(G\cap N_\delta(\ell))\geq\delta^{\eps-1}.
\]
\item For every $H\subset F\times G\times G$ with
$\cE_\delta(H)\geq\delta^{\eta-3}$, there are $\rho\geq\delta$
and an interval $I$ with $|I|\geq\delta^{-\eta}\rho$ for which
\[
 \cE_\rho\bigl(I\cap
       \{a\cdot(b_1-b_2):(a,b_1,b_2)\in H\}\bigr)
                  \geq(|I|/\rho)^{1-\eps}.
\]
\end{enumerate}
\end{theorem}
\begin{proof}
\uses{lem:hypergraph-refinement,prop:wz1-separated,
  lem:strip-localization}
Discretize the three vertex sets and refine the incidence graph so
that all its projections have comparable degrees. Localize the active
vertex sets into separated balls and apply the separated projection
alternative. If an endpoint family concentrates in a strip, rescale
that strip and apply the localized alternative. Transport the two
possible conclusions back to the original dot-difference expression.
The losses in the graph density and the scale are absorbed by choosing
$\eta$ after $\eps$.
\end{proof}

\begin{remark}[Which line passes through the origin]
\label{rem:projection-orientation}
WZ1 Theorem 5.2 places $F$ near the affine line and $G$ near the
perpendicular line through the origin. For the expression
$a\cdot(b_1-b_2)$ the roles are reversed: if $G$ lies on an affine
line with direction $v$, the degenerate set of coefficients is
$v^\perp$, which passes through the origin. The Lean statements use
this latter orientation, as does WZ1 Lemma 8.13.
For a concrete illustration, take
$F=\{0\}\times[1/2,3/4]$ and
$G=[1/2,3/4]\times\{1\}$. All the dot differences vanish, while
$G$ lies far from the horizontal line through the origin. Thus the
printed orientation cannot describe this degenerate example.
\end{remark}

\begin{remark}[Finite graphs and the two endpoint classes]
\label{rem:projection-model}
The linked Lean formulation uses separated finite sets, normalized
Frostman counts, and a uniformly dense tripartite graph. In the
application to the slope, the two endpoint classes are refinements of
one ambient set $G$. Some intermediate declarations therefore return
line concentration in at least one of those two classes, which is
enough to give concentration in $G$. This conclusion should not be
reported as simultaneous line concentration in two arbitrary,
independent endpoint sets; see Remark~\ref{rem:endpoint-union}.
\end{remark}

\section{Affine approximation on the retained heights}

\begin{lemma}[Affine approximation in one grain; WZ1, Lemma 5.3]
\label{lem:locally-affine}
\lean{Kakeya.Assouad.WZ1LocallyLinearOneScaleData}
\uses{def:ad-set,def:plane-map}
For every $\eps>0$, choose $\eta>0$ and a sufficiently small
$\rho>0$. Let $I=[z_*,z_*+\sqrt\rho]\subset[-1,1]$,
$Z\subset I$, and let $f:I\to[0,1]$ be $1$-Lipschitz. Suppose
$E\subset[-1,1]^2\times Z$ is a union of $\rho$-cubes lying within
$\sqrt\rho$ of a line in direction $(-f(z_*),1,0)$, with
\[
                       |E|\geq\rho^{1+\sigma/2+\eta}.
\]
Let $V:E\to S^2$ be $1$-Lipschitz with $|V(p)\cdot(0,0,1)|\leq1/2$.
Assume each horizontal projection
$(E\cap\{z=z_0\})\cdot(1,f(z_0),0)$ is a
$(\rho,1-\sigma,\rho^{-\eta})_1$-ADset, and
\[
       |(E\cap B(p,\sqrt\rho))\cdot V(p)|
                 \leq\rho^{(1+\sigma)/2-\eta}\quad(p\in E).
\]
Then there is $Z_{\rm lin}\subset Z$ of measure at least
$\rho^{1/2+\eps}$ and an affine function $L:I\to[0,1]$ satisfying
$|f(z)-L(z)|\leq\rho$ on $Z_{\rm lin}$.
\end{lemma}
\begin{proof}
\uses{thm:wz1-projection}
Choose well-populated horizontal slices and encode the local and
global normal directions by two planar graphs. The common shaded
incidences give a dense tripartite graph. The projection alternative
with a long, nearly one-dimensional dot-difference set would contradict
the global AD exponent $1-\sigma$. The remaining alternative forces
the slope graph to remain near an affine line on a substantial set
of heights. The volume assumption supplies the stated measure of
those heights.
\end{proof}

\begin{lemma}[Affine approximation at one scale; WZ1, Lemma 5.4]
\label{lem:piecewise-affine}
\lean{Kakeya.Assouad.WZ1LocallyLinearOneScaleConclusion}
\uses{prop:lipschitz-grains}
For every $\eps>0$, a sufficiently extremal grain configuration and
$\rho\in[\delta^{1-\eps},\delta^\eps]$ admit an
$\eps$-extremal refinement and a $2$-Lipschitz function
$L:[-1,1]\to\R$ which is affine on the grid intervals
$[m\sqrt\rho,(m+1)\sqrt\rho]$ and obeys
\[
        |f(z)-L(z)|\leq\rho
        \quad\text{for every retained height }z.
\]
\end{lemma}
\begin{proof}
\uses{lem:locally-affine,prop:sticky-structure}
Apply the one-grain result to a balanced decomposition of the source
configuration. Select compatible affine pieces in the height intervals
and retain the incidences above their good heights. Extend the chosen
pieces across the gaps with controlled slope. The balanced cover
provides the mass needed to restore extremality.
\end{proof}

\begin{lemma}[A finite affine hierarchy; WZ1, Lemma 5.5]
\label{lem:affine-hierarchy}
\lean{Kakeya.Assouad.WZ1LocallyLinearHierarchyData}
\uses{prop:lipschitz-grains}
Given $\eps>0$ and $N\geq1$, choose the source loss and scale
sufficiently small, and put $\rho_j=\delta^{j/N}$.
There is one $\eps$-extremal refinement with functions
$L_1,\ldots,L_N$, each $2$-Lipschitz and affine on intervals of length
$\sqrt{\rho_j}$, such that
\[
 |f(z)-L_j(z)|\leq\rho_j
 \quad(1\leq j\leq N, z\text{ a retained height}).
\]
\end{lemma}
\begin{proof}
\uses{lem:piecewise-affine}
Iterate at the prescribed scales, scheduling each input loss after
the next requested output loss. Subsequent refinements preserve every
previous approximation on the heights that survive.
\end{proof}

\begin{definition}[Vertical trapezoids; WZ1, Definition 5.1]
\label{def:trapezoid}
\lean{Kakeya.Assouad.WZ1VerticalTrapezoid}
\leanok
A vertical trapezoid has the form
\[
        R=\{(z,t):z\in J,\ |t-L(z)|\leq h\},
\]
where $J\subset[-1,1]$ is an interval and $L$ is affine. Its length
is $|J|$, its height parameter is $h$, and its slope is the slope
of $L$.
\end{definition}

\begin{corollary}[Nested trapezoids; WZ1, Corollary 5.6]
\label{cor:trapezoid-hierarchy}
\lean{Kakeya.Assouad.WZ1NestedSlopeCorrectionData}
\uses{def:trapezoid,prop:lipschitz-grains}
Given $\eps>0$ and $N\geq1$, one can refine a sufficiently extremal
grain configuration and find families $\mathcal R_j$ of trapezoids,
$\rho_j=\delta^{j/N}$, with the following properties. Every slope
lies in $[-2,2]$. At level $j$ the height is $\rho_j$ and the lengths
lie in $[\rho_j^{1/2+\eps},\sqrt{\rho_j}]$; the horizontal intervals
are $\sqrt{\rho_j}$-separated. Each trapezoid at level $j>1$ is
contained in a unique parent at level $j-1$. If $D$ is the set of
retained heights, then $\graph(f|_D)\subset\bigcup\mathcal R_N$.
\end{corollary}
\begin{proof}
\uses{lem:affine-hierarchy}
Select separated subintervals where the approximations at all levels
are compatible. Group their affine graphs into vertical trapezoids.
Refining the occupied height intervals at each step ensures actual
containment in the preceding trapezoids, as well as the mass retention
needed for extremality.
\end{proof}

\section{Smoothing the hierarchy}

\begin{lemma}[A smooth extension of a segment; WZ1, Lemma 5.8]
\label{lem:segment-extension}
\lean{Kakeya.Assouad.WZ1SmoothSegmentExtensionStatement,Kakeya.Assouad.wz1_smooth_segment_extension}
\leanok
Let $x_1<x_2$, let $L$ join $(x_1,y_1)$ to $(x_2,y_2)$, and let
$a>0$. There is a smooth function $G:\R\to\R$ agreeing with the
segment on $[x_1,x_2]$, vanishing outside
$[x_1-a,x_2+a]$, and satisfying
\[
 |G'|\lesssim h/a+|s|,\qquad
 |G''|\lesssim h/a^2+|s|/a,
 \quad h=|y_1|+|y_2|,\quad s=(y_2-y_1)/(x_2-x_1).
\]
\end{lemma}
\begin{proof}
\leanok
Multiply the affine function through the segment by smooth cutoffs
which equal one on the segment and turn off over intervals of length
$a$. Differentiating the product gives the two bounds. A single
absolute cutoff constant works for every segment and every $a$.
\end{proof}

\begin{proposition}[Smoothing nested trapezoids; WZ1, Proposition 5.7]
\label{prop:trapezoid-smoothing}
\lean{Kakeya.Assouad.WZ1MultiscaleSlopeCorrectionData}
\uses{def:trapezoid}
Let $N\geq2$, $\rho_j=\delta^{j/N}$, and suppose the nested
trapezoid families have slopes in $[-2,2]$, heights $\rho_j$,
lengths in $[\rho_j^{1/2+\eta},\sqrt{\rho_j}]$, and horizontal
separation $\sqrt{\rho_j}$ at level $j$. If
$\graph(f)\subset\bigcup\mathcal R_N$, $f:D\to\R$,
$D\subset[-1,1]$, then there is $g\in C^2([-1,1])$ satisfying
\[
 |f-g|\leq\delta\text{ on }D,
 \qquad |g''|\lesssim\delta^{-1/N-2\eta}.
\]
On the horizontal interval of $R\in\mathcal R_j$, the derivative
differs from the slope of $R$ by at most a constant times
$\rho_j/\rho_{j+1}^{1/2+\eta}$, for $j<N$.
\end{proposition}
\begin{proof}
\uses{lem:segment-extension}
Begin with the coarsest affine pieces. At each finer level, extend
the difference between a child's affine function and its parent's
using the segment-extension lemma. Their supports are separated at
that level. The resulting finite sum agrees with the terminal affine
piece wherever the fine graph is retained. Summing the first- and
second-derivative costs over levels gives the estimates.
\end{proof}

\begin{remark}[Constants and terminal indices in smoothing]
\label{rem:smoothing}
The paper's derivative comparison uses $\rho_{j+1}$; with only
$\rho_1,\ldots,\rho_N$ defined, that formula is read for $j<N$.
The Lean construction records the first- and second-derivative cost
of every supported correction and sums them explicitly, including
the terminal correction. It also records disjointness of supports
and compatibility of parent and child intervals. These conditions
are needed before a bound for each individual correction can be
turned into a bound for their sum.
\end{remark}
